% Source: https://artofproblemsolving.com/community/c3h1353162_counting_cubes_in_cubes
% Stable identity: aops-topic-1353162-q01
Use coordinates with $O=(0,0,0)$, $Q=(5,5,5)$, the $z$-axis vertical. The two vertical edges in question are $x=5,y=0$ and $x=0,y=5$, with midpoints $P=(5,0,2.5)$ and $R=(0,5,2.5)$. All four points satisfy
\[
x+y=2z ,
\]
which is therefore the plane; it cuts the cube in the rhombus $OPQR$.

\begin{asy}
usepackage("amsmath");
size(5.5cm);
defaultpen(fontsize(10pt));
pair pr(triple p){return (p.x+0.45*p.y, p.z+0.3*p.y);}
triple[] V={(0,0,0),(5,0,0),(5,5,0),(0,5,0),(0,0,5),(5,0,5),(5,5,5),(0,5,5)};
pair v(int i){return pr(V[i]);}
pair O=pr((0,0,0)),P=pr((5,0,2.5)),Q=pr((5,5,5)),R=pr((0,5,2.5));
fill(O--P--Q--R--cycle,paleblue+opacity(0.6));
draw(v(0)--v(1)--v(2),linewidth(0.9)); draw(v(4)--v(5)--v(6)--v(7)--cycle,linewidth(0.9));
draw(v(0)--v(4),linewidth(0.9)); draw(v(1)--v(5),linewidth(0.9)); draw(v(2)--v(6),linewidth(0.9));
draw(v(0)--v(3),gray+dashed); draw(v(3)--v(2),gray+dashed); draw(v(3)--v(7),gray+dashed);
draw(O--P--Q,blue); draw(Q--R--O,blue+dashed);
dot(O);dot(P);dot(Q);dot(R);
label("$O$",O,SW); label("$Q$",Q,NE); label("$P$",P,E); label("$R$",R,W);
label("\mbox{}",(min(currentpicture,user=true)+max(currentpicture,user=true))/2,black); // keeps the text after the figure black when LaTeX typesets the labels inline
\end{asy}

The unit cube $[i,i+1]\times[j,j+1]\times[k,k+1]$ ($0\le i,j,k\le4$) has interior points on both sides of the plane exactly when the expression $x+y-2z$ takes both signs on its interior. On the cube this expression ranges over the open interval from $i+j-2k-2$ to $i+j-2k+2$ in the interior, so the condition is $|i+j-2k|\le1$, i.e.\ $i+j\in\{2k-1,2k,2k+1\}$.

The number of pairs $(i,j)$ with $i+j=t$ is $1,2,3,4,5,4,3,2,1$ for $t=0,\dots,8$. Layer by layer:
\[
k=0:\ 0+1+2=3,\quad k=1:\ 2+3+4=9,\quad k=2:\ 4+5+4=13,\quad k=3:\ 4+3+2=9,\quad k=4:\ 2+1+0=3 .
\]
The plane cuts $3+9+13+9+3=\boxed{37}$ unit cubes.
